\section{Repository Architecture}

The repository is organized around a scientific preprocessing and physics stack
for a microfluidic asymmetric-loop device. The architecture now follows the
physical workflow:

\[
\mathrm{Geometry}
\rightarrow
\mathrm{Occupancy}
\rightarrow
\mathrm{Hydraulics}
\rightarrow
\mathrm{CFD}
\rightarrow
\mathrm{Physics~State}.
\]

The code implementing this workflow lives under:

\begin{verbatim}
src/physics/
  geometry/
  occupancy/
  hydraulics/
  cfd/
  state/
\end{verbatim}

The current implementation is deterministic and physics-preparatory. It does
not train a transformer, optimize temporal lag, fit hydraulic parameters, or
resolve droplets as moving multiphase CFD objects. It does include a validated
steady Stokes full-device CFD solve for the continuous phase.

\subsection{Top-Level Layout}

\begin{longtable}{@{}ll@{}}
\toprule
Path & Scientific responsibility \\
\midrule
\texttt{src/config/} & Resolve experiment and device configuration. \\
\texttt{src/physics/geometry/} & Device centerlines, physical geometry, and region labels. \\
\texttt{src/physics/occupancy/} & Droplet fractional occupancy over physical regions. \\
\texttt{src/physics/hydraulics/} & Baseline branch-resistance and superficial-velocity model. \\
\texttt{src/physics/full\_device\_cfd/} & Full-loop CAD geometry, mesh generation, Stokes solve, and field verification. \\
\texttt{src/physics/state/} & Placeholder for future unified physics state. \\
\texttt{scripts/geometry/} & Geometry and device-region command-line tools. \\
\texttt{scripts/occupancy/} & Droplet occupancy command-line tools. \\
\texttt{scripts/physics/} & Hydraulics and CFD command-line tools. \\
\texttt{scripts/validation/} & Standalone scientific validation analyses. \\
\bottomrule
\end{longtable}

\subsection{Physical Device}

The present device is an asymmetric microfluidic loop with channel width and
height both set to \(100~\mu\mathrm{m}\). The image calibration is:

\[
1~\mathrm{px} = 4~\mu\mathrm{m}.
\]

The loop branch lengths are:

\[
L_{\mathrm{left}} = 1791~\mu\mathrm{m},
\qquad
L_{\mathrm{right}} = 1491~\mu\mathrm{m}.
\]

Thus the left branch is the long branch and the right branch is the short
branch:

\[
\frac{L_{\mathrm{left}}}{L_{\mathrm{right}}} = 1.201207,
\qquad
\alpha = \frac{L_{\mathrm{left}} - L_{\mathrm{right}}}{L_{\mathrm{right}}}
= 0.201207.
\]
